% PAQUETES

% Paquetes de idioma y codificación principales
\usepackage[utf8]{inputenc} % Permite usar caracteres UTF-8.
\usepackage[T1]{fontenc} % Mejora la codificación de las fuentes.
\usepackage[spanish,es-nodecimaldot]{babel} % Adapta el documento al idioma español.
\usepackage{csquotes} % Gestiona correctamente las comillas.
\usepackage{graphicx} % Permite insertar y ajustar imágenes.
\usepackage{booktabs} % Crea tablas con formato profesional.
\usepackage{amsmath} % Añade herramientas para fórmulas matemáticas.
\usepackage{lmodern} % Usa la familia tipográfica Latin Modern.
\usepackage{caption} % Personaliza el formato de títulos de figuras y tablas.
\usepackage[hidelinks]{hyperref} % Enlaces PDF sin bordes ni colores.
\usepackage{bookmark} % Mejora la generación de marcadores PDF.
% Gestiona citas y bibliografía con Biber.
\usepackage[
    backend=biber,
    style=authoryear,
    language=spanish,
    sorting=nyt,
    maxcitenames=2
]{biblatex}

% Formato autor-año con coma, equivalente al uso habitual de APA.
\DeclareDelimFormat{nameyeardelim}{\addcomma\space}

\addbibresource{bibliography/references.bib} % Gestiona las referencias bibliográficas en formato APA.

% Paquetes de diseño y formato
\usepackage{titlesec}
\usepackage{xcolor} % Permite definir y usar colores personalizados.
\usepackage{tikz} % Permite crear gráficos y formas vectoriales.
\usepackage{tocloft} % Personaliza la tabla de contenido.
%\usepackage{layout} % Muestra las dimensiones y márgenes junto con sus propiedades de la página con \layout.
%\usepackage{showframe} % Visualizar los bordes del cuerpo y sus dimensiones.
% Configura los márgenes y el área útil de la página.
\usepackage[top=3cm,bottom=3cm,left=4cm,right=2cm]{geometry}